% ************************************************************
% ABSTRACT
% ************************************************************

\begin{abstract}
This thesis develops a reproducible computational pipeline for evaluating whether persistent-homology summaries can serve as auxiliary diagnostics for structure in symmetric-cipher output and generated byte-stream controls. The study compares AES-128 CTR deterministic-output conditions, a seeded SHA-256 expansion reproducibility baseline, operating-system cryptographic random output as a non-deterministic sensitivity condition, weak-generator controls, and deliberately structured periodic and biased byte-stream controls. Streams are generated locally from documented keys, nonces, plaintext patterns, seeds, replicate indices, generator parameters, and software versions, avoiding dependence on restricted datasets. The registered analyses use byte-pair and sliding-window embeddings, H0/H1 persistent-homology summaries, internal randomness diagnostics, effect-size tables, validation gates, and evidence registers for claim traceability. Across the current evidence base, deliberately structured weak controls separate strongly from the SHA-256 seeded baseline under the registered topological summaries, and the weak-control sensitivity pattern reproduces under an alternate deterministic master seed. AES-CTR, OS CSPRNG, and xorshift32 do not show comparable top-ranked separation under the tested design. The thesis does not propose a cryptanalytic attack, identify a cipher vulnerability, or claim to certify cryptographic security. Its contribution is a transparent, object-oriented, reproducible framework for testing whether topological data analysis can supplement established statistical diagnostics in cryptographic implementation review.
\end{abstract}
